\documentclass[11pt]{article}
\pdfinfoomitdate=1
\pdfsuppressptexinfo=15
\usepackage[T1]{fontenc}
\usepackage{lmodern}
\usepackage[margin=1in]{geometry}
\usepackage{microtype}
\usepackage{amsmath,amssymb,amsthm,mathtools}
\usepackage{enumitem,needspace}
\usepackage{tikz-cd}
\usepackage[dvipsnames]{xcolor}
\PassOptionsToPackage{hyphens}{url}
\usepackage{hyperref}
\usepackage{bookmark}
\hypersetup{colorlinks=true,linkcolor=MidnightBlue,citecolor=MidnightBlue,
  urlcolor=MidnightBlue,pdftitle={A universal group of type F infinity},
  pdfauthor={OpenAI},pdfsubject={Universal groups of geometric type F infinity}}
\newcommand{\Z}{\mathbb Z}
\newcommand{\N}{\mathbb N}
\newcommand{\F}{\mathbb F}
\newcommand{\Finfty}{F_\infty}
\newcommand{\FPinfty}{FP_\infty}
\newtheorem{theorem}{Theorem}[section]
\newtheorem{lemma}[theorem]{Lemma}
\newtheorem{proposition}[theorem]{Proposition}
\newtheorem{corollary}[theorem]{Corollary}
\theoremstyle{definition}
\newtheorem{definition}[theorem]{Definition}
\theoremstyle{remark}
\newtheorem{remark}[theorem]{Remark}
\numberwithin{equation}{section}
\setlist[enumerate]{topsep=5pt,itemsep=3pt}
\setlength{\emergencystretch}{2em}
\title{A universal group of type $\Finfty$}
\author{OpenAI}
\date{September 23, 2026}

\pdfgentounicode=1
\pdfglyphtounicode{parenleftbig}{0028}
\pdfglyphtounicode{parenrightbig}{0029}
\pdfglyphtounicode{angbracketleftbig}{27E8}
\pdfglyphtounicode{angbracketrightbig}{27E9}
\pdfglyphtounicode{parenlefttp}{239B}
\pdfglyphtounicode{parenrighttp}{239E}
\pdfglyphtounicode{parenleftbt}{239D}
\pdfglyphtounicode{parenrightbt}{23A0}

\begin{document}
\maketitle
\begin{abstract}
We construct a single group of type $\Finfty$ containing every finitely
presented group. Its finitely generated subgroups, up to isomorphism, are
exactly the finitely generated recursively presented groups. This answers
the $\Finfty$ form of the higher-dimensional Higman embedding question.
\end{abstract}

\section{Introduction}\label{sec:introduction}

A group has type $F_n$ if it admits a classifying CW complex with finite
$n$-skeleton, and type $\Finfty$ if it admits one classifying complex
with finitely many cells in every dimension. Here a classifying complex
is a connected CW complex with the given fundamental group and
contractible universal cover; it may have infinitely many dimensions.
Type $F_2$ is equivalent to finite presentability.

\begin{theorem}\label{thm:main}
There exists a group $H$ of type $\Finfty$ such that every finitely
presented group admits an injective homomorphism into $H$.
\end{theorem}

The common target gives an exact description of its finitely generated
subgroups. A finitely generated group is recursively presented if it
admits a presentation on finitely many generators with a recursively
enumerable set of relators.

\begin{corollary}\label{cor:recursive}
A finitely generated group embeds in $H$ if and only if it is
recursively presented.
\end{corollary}

\begin{proof}
If $G$ is finitely generated and recursively presented, Higman's
embedding theorem places it in a finitely presented
group~\cite{Higman1961}, which embeds in $H$ by
Theorem~\ref{thm:main}. Conversely, $H$ is finitely presented because
it has type $\Finfty$. For a finitely generated subgroup $G\le H$,
represent its generators by words in a finite presentation of $H$.
Enumerate all words in the chosen generators, and dovetail finite
derivations from the ambient relators proving that their substituted
words are trivial. Recording the successful words enumerates exactly
the defining relators of $G$.
\end{proof}

This answers the $\Finfty$ form of the higher-dimensional Higman
embedding question, stated for finitely generated recursively presented
groups in~\cite[Question~1.3]{FFZ2026}. The theorem imposes no
word-problem hypothesis on the input.

\subsection*{Context and construction}

For a commutative unital coefficient ring $\Lambda$, a group $G$ has
type $FP_n(\Lambda)$ if the trivial $\Lambda G$-module $\Lambda$ has a
projective resolution with finitely generated terms through degree $n$;
write $FP_n$ for $\Lambda=\Z$. For $n=\infty$, finite generation is
required in every degree. Leary proved that every countable group embeds
in a group of type $FP_2$~\cite{Leary2018}. Fournier-Facio and Zaremsky
show that, for each $n\in\N\cup\{\infty\}$, embeddings of all finitely
generated recursively presented groups into recursively presented
groups of type $FP_n$ would imply embeddings of all finitely presented
groups into groups of type $F_n$~\cite[Theorem~A]{FFZ2026}.
They also prove that the rope-trick groups of their Definition~1.5 fail
type $FP_3(\mathbb Q)$ for infinite inputs with nontrivial
free-presentation kernel~\cite[Definition~1.5 and Theorem~B]{FFZ2026}.
This obstruction is specific to those rope-trick groups.

The proof of their conditional theorem uses a universal finitely
presented group, an embedding through an intermediate overgroup and
back, and an ascending HNN extension. They relate this strategy to
earlier work of Baumslag--Dyer--Miller and
Fournier-Facio--L\"oh--Moraschini~\cite[Section~2]{FFZ2026}.

For the universal finitely presented group $U$ of
Lemma~\ref{lem:universal}, Proposition~\ref{prop:finite-envelope}
gives a finite presentation of a group $\widehat W$ equipped with two
finite families of homomorphisms from $U$, sharing a distinguished map
$\iota:U\to\widehat W$. Faithful models in tensor product algebras
verify that $\iota$ is injective. Finite presentability of $\widehat W$ then lets
universality supply an embedding $\kappa:\widehat W\to U$:
\[
\begin{tikzcd}[column sep=large]
 U \arrow[r,hook,"\iota"]
   \arrow[rr,bend right=24,hook,"f=\kappa\iota"']
 & \widehat W \arrow[r,hook,"\kappa"] & U.
\end{tikzcd}
\]
Postcomposing the two families by $\kappa$ gives homomorphisms from $U$
to itself, with the injective endomorphism $f$ as their common
distinguished map (Corollary~\ref{prop:self-map}). The group $U$ embeds
in the associated ascending HNN extension
\[
 H=T(U,f)=\langle U,t\mid t^{-1}ut=f(u)\ (u\in U)\rangle,
\]
which is finitely presented by Lemma~\ref{lem:torus-presentation}.
Theorem~\ref{thm:torus} proves that $H$ has type $\Finfty$.

The finite idempotent systems, the group enlargement, and the
ascending-torus criterion are shared with the
companion~\cite[Sections~2.1, 2.3--2.5, and~3.1]{OpenAIOvergroups}.
We give local proofs of the algebraic identities and geometric
criterion.

The geometric step begins with a uniform filling statement for ordinary
integral chains. In the homogeneous chain model, restricting relative
vertex differences to a finite subset of $U$ leaves only finitely many
simplex types up to common left translation in each degree. For each
positive degree and finite input set of relative differences,
Proposition~\ref{prop:uniform-fillings} chooses one exponent and one
finite output set. That iterate of $f$ sends every integral cycle using
the input differences to the boundary of a chain using the output
differences. The filling lengths and coefficients may vary with the cycle.

The finite diagrams supply these fillings by induction. For each
$\ell\in\{2,3\}$, the diagram is indexed by idempotents over
$\F_\ell$. Orthogonal sums correspond to pointwise products of
homomorphisms with elementwise commuting images.
Lemma~\ref{lem:idempotents} says that an additive assignment has its
value at the identity annihilated by $\ell$. In degrees $n\ge2$, the
chain diagonal includes mixed terms from lower degrees, so the
corresponding cycle classes are not automatically additive. Using
lower-degree fillings, the induction constructs, below degree $n$, a
controlled chain homotopy from the map induced by an iterate of $f$ to
the constant map. Lemma~\ref{lem:controlled-additivity} uses it to
remove the mixed terms. For each input cycle, both diagrams then give
additive assignments into the homology of the same bounded chain
complex. Their common value at the identity idempotent is represented
by the image of the input cycle under an iterate of $f$. It is
annihilated by both $2$ and $3$, and hence vanishes.

For an arbitrary set $\mathcal I$, put $M=(\Z H)^{\mathcal I}$. A cycle
in the tensor complex computing $H_i(U,M)$, $i>0$, decomposes in each
coordinate into finitely many ordinary cycles in left $U$-cosets, with
one finite set of simplex types common to all coordinates. The uniform
filling statement supplies one exponent and output set for these cycles,
so their fillings reassemble into one tensor chain. In the ascending HNN
sequence of Section~\ref{sec:products}, the map $F$ combines $f$ on
chains with multiplication by the stable letter on coefficients. The
fillings make $F$ locally nilpotent on $H_i(U,M)$ for $i>0$: every class
is killed by some power of $F$. Thus $1-F$ is invertible there, and
exactness gives $H_n(H,M)=0$ for every $n\ge2$. Since $H$ is finitely
presented, Lemma~\ref{lem:product-criterion} turns this vanishing into
geometric finiteness. At each stage, the cellular cycle module is
finitely generated; via Hurewicz, those generators have sphere
representatives. Attaching cells along these spheres stage by stage
produces one classifying complex with finitely many cells in every
dimension.

\subsection*{Conventions}

All groups are discrete. Group homology has integral coefficients
unless stated otherwise, and coefficient modules are right modules.
Products indexed by arbitrary sets always mean direct products, not
direct sums.

\section{A universal finitely presented group}
\label{sec:universal}

We will place a finite collection of homomorphisms inside one finitely
presented overgroup and then embed that overgroup back into the original
group.  For this purpose we need a finitely presented group containing
every finitely presented group.  We recall its construction from
Higman's embedding theorem \cite{Higman1961}: every finitely generated
group with a recursively enumerable presentation embeds in a finitely
presented group.  No decidability assumption on the word problem occurs
in this theorem.

The base-group embedding and finite-generation construction below
follow the classical embedding methods of Higman, B.~H.~Neumann and
Hanna Neumann~\cite[Theorems~I, I', and~IV]{HNN1949}; we include the arguments
in the form needed here. We first record an elementary form of
base-group injectivity for HNN extensions. Here the countable case is
sufficient.

\begin{lemma}\label{lem:hnn-injection}
Let $D$ be a countable group, let $D_1,D_2$ be subgroups of $D$, and let
$\theta:D_1\to D_2$ be an isomorphism.  The natural homomorphism from $D$
to
\[
 \bigl\langle D,t\mid t^{-1}dt=\theta(d)\quad(d\in D_1)\bigr\rangle
\]
is injective.  Here the presentation includes all relations of $D$.
\end{lemma}

\begin{proof}
Let $D$ act faithfully on $\Omega=D\times\N$ by left multiplication
in the first coordinate.  The restrictions to $D_1$ and $D_2$ are free
actions, each with countably infinitely many orbits.  Match their orbit
sets bijectively.  For each matched pair choose a base point $x$ in the
$D_2$-orbit and a base point $y$ in the $D_1$-orbit, and define
\[
 T\bigl(\theta(d)x\bigr)=dy\qquad(d\in D_1).
\]
Freeness and the orbit matching make $T$ a permutation of $\Omega$.
For every $d\in D_1$ it satisfies $T\theta(d)=dT$, and hence
$T^{-1}dT=\theta(d)$, as permutations of $\Omega$.  Sending $t$ to $T$
therefore defines an action of the displayed group whose restriction to
$D$ is faithful.  The natural map from $D$ must be injective.
\end{proof}

\begin{lemma}\label{lem:universal}
There exists a finitely presented group $U$ into which every finitely
presented group embeds.
\end{lemma}

\begin{proof}
Effectively list all finite presentations, renaming their generators
so that the generating alphabets are disjoint.  Let $C$ be the group
given by the union of these presentations.  Its generators can be
effectively numbered $c_0,c_1,\ldots$, and its relators can be
effectively enumerated.  Each group on the list embeds in $C$: the
homomorphism induced by its generators has a left inverse obtained by
killing all generators of the other presentations.

It remains to place $C$ inside a finitely generated group with
recursively enumerable relators.  Form
\[
 D=C*\langle a,p\rangle,
 \qquad v_i=a^ipa^{-i}\quad(i\geq0),
\]
where $\langle a,p\rangle$ is free of rank two.  The elements $v_i$
freely generate their subgroup: in any product of nonzero powers of
the $v_i$ with consecutive indices distinct, consecutive powers of
$p$ are separated by nonzero powers of $a$, so free reduction cannot
give the identity.  The elements $c_iv_i$ also freely generate their
subgroup, because the projection $D\to\langle a,p\rangle$ killing $C$
sends them to the $v_i$.  Consequently $v_i\mapsto c_iv_i$ extends to
an isomorphism between these two free subgroups of $D$.

By Lemma~\ref{lem:hnn-injection}, $D$, and hence $C$, embeds in
\[
 V=\bigl\langle D,q\mid q^{-1}v_iq=c_iv_i\quad(i\geq0)\bigr\rangle.
\]
The identities
\[
 c_i=q^{-1}v_iqv_i^{-1}
\]
show that $V$ is generated by $a,p,q$.  More precisely, it has a
presentation on these three generators obtained by substituting the
displayed word for every $c_i$ in the relators of $C$.  The substitution
turns each HNN relation into an identity.  Thus it and the natural map
into $V$ are mutually inverse homomorphisms, so this is indeed a
presentation of $V$.  Its relators are recursively
enumerable.  Higman's embedding theorem now embeds $V$ in a finitely
presented group $U$.  Every finitely presented group embeds in $C$,
then in $V$, and finally in $U$.
\end{proof}

Below we use only the existence of the subgroup embeddings supplied by
Lemma~\ref{lem:universal}.

\section{Two finite systems of idempotents}
\label{sec:idempotents}

The finite diagrams used in our construction are indexed by idempotent
operators over $\F_2$ and $\F_3$.  Their essential property is elementary:
an additive assignment to either diagram forces the value at the identity
to be annihilated by the corresponding prime.  These diagrams and their
additivity argument are shared with the
companion~\cite[Section~2.1, Lemma~2.1]{OpenAIOvergroups}.
We give the complete argument before realizing them by homomorphisms of
groups.

For $\ell\in\{2,3\}$, let $D_\ell$ be the unital $\F_\ell$-algebra of
operators on the free vector space with basis $\{0,1\}^{\N}$ generated by
$l_0,l_1,r_0,r_1$.  Here $l_i$ prefixes the symbol $i$, and $r_i$ deletes
an initial $i$, sending basis vectors with the other initial symbol to
zero.  Writing $I$ for the identity operator, we have
\begin{equation}\label{eq:binary-prefix-relations}
 r_i l_j=\delta_{ij}I,
 \qquad l_0r_0+l_1r_1=I.
\end{equation}
For a finite binary word $v$, let $l_v$ prefix $v$, let $r_v$ delete $v$
when it occurs as a prefix and give zero otherwise, and put $E_v=l_vr_v$.
Thus $E_v$ projects onto the span of the sequences beginning with $v$.

If $(v_1,v_2,v_3)$ is a prefix-free list, then
\[
 (l_{v_i}r_{v_j})(l_{v_k}r_{v_m})
   =\delta_{jk}l_{v_i}r_{v_m}.
\]
These operators are nonzero matrix units, so their span is a copy of
$M_3(\F_\ell)$.  We use three such algebras, connected by shared
diagonal projections.  Write
\[
 A=E_{00},\qquad A_0=E_{000},\qquad A_1=E_{001},
 \qquad B'=E_{01},\qquad C'=E_1.
\]
Their defining lists, diagonal projections, and identity operators are
\begin{equation}\label{eq:three-prefix-lists}
\begin{array}{c|c|c}
 \text{prefix list} & \text{diagonal projections} & \text{identity}\\ \hline
 (00,01,1)  & A,B',C'   & I\\
 (000,01,1) & A_0,B',C' & A_0+B'+C'\\
 (001,01,1) & A_1,B',C' & A_1+B'+C'.
\end{array}
\end{equation}
The last two identities are not $I$, while
$A=A_0+A_1$ and $A_0A_1=A_1A_0=0$.  Let $\mathcal P_\ell$ be the union,
as a set of actual
operators in $D_\ell$, of all idempotents in these three matrix
algebras.  It is finite and contains $I$.  An operator belonging to two
of the algebras has just one index in $\mathcal P_\ell$.

Two operators $E,P$ are called \emph{orthogonal} if $EP=PE=0$.
The relations below include all orthogonal pairs whose sum belongs to
$\mathcal P_\ell$, even when the pair does not belong to a single one of
the three matrix algebras.  In particular, they include
$A=A_0+A_1$, whose summands come from the two smaller support algebras.

\begin{lemma}\label{lem:idempotents}
Let $\ell\in\{2,3\}$, let $T$ be an abelian group, and let
$(z_E)_{E\in\mathcal P_\ell}$ be elements of $T$.  Suppose that
\begin{equation}\label{eq:idempotent-additivity}
 z_{E+P}=z_E+z_P
\end{equation}
whenever $E,P,E+P\in\mathcal P_\ell$ and $EP=PE=0$.
Then $\ell z_I=0$.
\end{lemma}

\begin{proof}
We first compare the values of rank-one idempotents within any one of
the three matrix algebras.  For a column vector $u\in\F_\ell^3$ and a row
covector $\lambda$ with $\lambda u=1$, write
\[
 x(u,\lambda)=z_{u\lambda}.
\]
Fix a nonzero vector $v\in\ker\lambda$.  We claim that the difference
\begin{equation}\label{eq:fixed-increment}
 x(u+v,\lambda)-x(u,\lambda)
\end{equation}
does not depend on the vector $u$ subject to $\lambda u=1$.

Take another such vector $u'$.  Suppose first that $u'-u$ and $v$ are
linearly independent.  They form a basis of $\ker\lambda$, and hence
$u,u'-u,v$ form a basis of $\F_\ell^3$.  There is consequently a covector
$\mu$ satisfying
\[
 \mu u=\mu u'=0,\qquad \mu v=1.
\]
The identity
\[
 u\lambda+v\mu=(u+v)\lambda+v(\mu-\lambda)
\]
expresses one rank-two idempotent as two orthogonal sums of rank-one
idempotents.  Indeed, $\lambda v=\mu u=0$ on the left, while
$\lambda v=(\mu-\lambda)(u+v)=0$ on the right; all four rank-one
factors have their defining covector equal to one on their defining
vector.  Applying \eqref{eq:idempotent-additivity} gives
\begin{equation}\label{eq:rank-one-difference}
 x(u+v,\lambda)-x(u,\lambda)
   =x(v,\mu)-x(v,\mu-\lambda).
\end{equation}
The same covector $\mu$ gives the same right-hand side for $u'$, proving
the claim in this case.

If $u'-u$ lies in the span of $v$, choose
$w\in\ker\lambda\setminus\langle v\rangle$ and put $u''=u+w$.
Both $u''-u$ and $u''-u'$ are independent of $v$.  Comparing first $u$
with $u''$ and then $u'$ with $u''$ proves the claim in general.
The kernel of $\lambda$ has dimension two, so this argument also works
over $\F_2$; no choice of a third nonzero scalar is being used.

Apply the claim to the successive vectors
$u,u+v,\ldots,u+(\ell-1)v$.  The $\ell$ successive differences are
equal, and they sum to zero because $\ell v=0$.  Therefore
\begin{equation}\label{eq:rank-one-prime-invariance}
 \ell x(u+v,\lambda)=\ell x(u,\lambda).
\end{equation}
The same conclusion is immediate for $v=0$, so $\ell x(u,\lambda)$
is independent of $u$ whenever $\lambda$ is fixed and $\lambda u=1$.
For the covector comparison, work within this same matrix algebra
and define a new assignment $z'_E=z_{E^{\mathsf T}}$.
Transposition preserves idempotence and orthogonal sums, so $z'$ obeys
the same additivity relations.  Apply the invariance just proved to
$x'(v,\mu)=z'_{v\mu}=z_{\mu^{\mathsf T}v^{\mathsf T}}$.
Fixing $\mu$ and varying $v$ now shows that $\ell x(u,\lambda)$
is independent of $\lambda$ whenever $u$ is fixed and $\lambda u=1$.
This does not assume that $z_E=z_{E^{\mathsf T}}$.

In particular the three diagonal matrix units have equal
$\ell$-multiples.  To see the comparison explicitly, for distinct basis
vectors $e_i,e_j$ take $\lambda=e_i^*+e_j^*$.  First change the covector
at $e_i$ from $e_i^*$ to $\lambda$, then change the vector from $e_i$
to $e_j$, and finally change the covector to $e_j^*$.  Each change
preserves the $\ell$-multiple, giving
\[
 \ell z_{e_i e_i^*}=\ell z_{e_j e_j^*}.
\]

It remains to use the shared projections displayed above.  The
comparison of diagonal units in each algebra gives
\[
 \ell z_A=\ell z_{B'}=\ell z_{C'}
           =\ell z_{A_0}=\ell z_{A_1}.
\]
Call this common value $q\in T$.  The projection identity
$A=A_0+A_1$ is an orthogonal sum across the two smaller support
algebras.  Its three operators belong to $\mathcal P_\ell$, so
\eqref{eq:idempotent-additivity} gives $q=q+q$.  Thus $q=0$.
Finally, $A+B'$ belongs to the first matrix algebra, and
$I=(A+B')+C'$.  Applying additivity twice yields
\[
 \ell z_I=\ell z_A+\ell z_{B'}+\ell z_{C'}=0,
\]
as required.
\end{proof}

\section{Faithful finite diagram envelopes}
\label{sec:envelopes}

The finite idempotent sets $\mathcal P_2$ and $\mathcal P_3$ specify
only finitely many commutation and multiplication identities.  We now
realize all these identities in a finitely presented overgroup, without
collapsing the common full-identity copy of the input group.  A faithful
model supplies the injectivity; its own presentation need not be finite.

\begin{proposition}\label{prop:finite-envelope}
For every finitely presented group $U$ there are a finitely presented
group $\widehat W$ and homomorphisms
\[
 k_{\ell,E}:U\longrightarrow\widehat W
 \qquad(\ell\in\{2,3\},\ E\in\mathcal P_\ell)
\]
such that $k_{2,I}=k_{3,I}$ is injective.  For each $\ell$, whenever
$E,E',E+E'\in\mathcal P_\ell$ and $EE'=E'E=0$, the subgroups
$k_{\ell,E}(U)$ and $k_{\ell,E'}(U)$ commute elementwise, and
\[
 k_{\ell,E+E'}(u)=k_{\ell,E}(u)k_{\ell,E'}(u)
 \qquad(u\in U).
\]
\end{proposition}

\begin{proof}
Fix a finite presentation $U=\langle X\mid\mathcal R\rangle$.
For every pair $(\ell,E)$ take a copy of this presentation, with
generators $x_{\ell,E}$ for $x\in X$.  Define $\widehat W$ by these
copy presentations and the following additional relations:
\begin{align}
 x_{2,I}&=x_{3,I} &&(x\in X),\label{eq:envelope-common}\\
 [x_{\ell,E},y_{\ell,E'}]&=1 &&(x,y\in X),\label{eq:envelope-commute}\\
 x_{\ell,E+E'}&=x_{\ell,E}x_{\ell,E'} &&(x\in X).
 \label{eq:envelope-sum}
\end{align}
The last two families run over both values of $\ell$ and every
orthogonal pair $E,E'$ whose members and sum lie in
$\mathcal P_\ell$.  All three lists are finite.  The indices are actual
idempotents in $D_\ell$, not separately tagged copies from its three
matrix subalgebras.  Thus these relations include the orthogonal sums
between different subalgebras as well as those within one subalgebra.
In particular, the zero idempotent is included; its sum relation with
itself gives $x_{\ell,0}=x_{\ell,0}^2$, forcing
$x_{\ell,0}=1$.

The copy presentations give homomorphisms $k_{\ell,E}:U\to\widehat W$.
Equation~\eqref{eq:envelope-commute} implies that their corresponding
image subgroups commute elementwise.  It follows that
$u\mapsto k_{\ell,E}(u)k_{\ell,E'}(u)$ is a homomorphism.  By
Equation~\eqref{eq:envelope-sum}, it agrees with $k_{\ell,E+E'}$ on
$X$, and therefore on all of $U$.  Similarly,
Equation~\eqref{eq:envelope-common} gives $k_{2,I}=k_{3,I}$.

To prove that this common map is injective, we construct a faithful
model.  The same group-level enlargement is used in the
companion~\cite[Section~3.1]{OpenAIOvergroups}; we verify all of its needed algebraic
properties here.  For any group $V$
and $\ell\in\{2,3\}$, work in the unital algebra
\[
 A_\ell(V)=D_\ell\otimes_{\F_\ell}\F_\ell[V],
\]
where $\F_\ell[V]$ is the group algebra with basis indexed by $V$.
For $E\in\mathcal P_\ell$ define
\begin{equation}\label{eq:envelope-model}
 j_{\ell,E}(v)=(I-E)\otimes1+E\otimes v
 \qquad(v\in V).
\end{equation}
Since $E$ and $I-E$ are orthogonal idempotents, multiplication gives
\[
 j_{\ell,E}(v)j_{\ell,E}(w)=j_{\ell,E}(vw),
 \qquad j_{\ell,E}(1)=1.
\]
Thus these elements are units, with inverses $j_{\ell,E}(v^{-1})$,
and $j_{\ell,E}$ is a homomorphism.

For $E,E',E+E'\in\mathcal P_\ell$ with $EE'=E'E=0$, rewrite
Equation~\eqref{eq:envelope-model} as $1+E\otimes(v-1)$.  Thus
\[
 \bigl(1+E\otimes(v-1)\bigr)
 \bigl(1+E'\otimes(w-1)\bigr)
 =1+E\otimes(v-1)+E'\otimes(w-1),
\]
because $EE'=0$.  Reversing the order gives the same answer since
$E'E=0$.  Hence the images of $j_{\ell,E}$ and $j_{\ell,E'}$
commute elementwise, even if $v$ and $w$ do not commute.  Setting
$w=v$ gives
\[
 j_{\ell,E}(v)j_{\ell,E'}(v)=j_{\ell,E+E'}(v).
\]
Moreover, $j_{\ell,I}(v)=I\otimes v$ is injective: the decomposition
$A_\ell(V)=\bigoplus_{v\in V}D_\ell\otimes v$ distinguishes these
elements, since $I\ne0$.

Let $W_\ell(V)$ be the subgroup of $A_\ell(V)^\times$ generated by
all images $j_{\ell,E}(V)$.  First form $W_2(U)$, and then form
$W_3(W_2(U))$.  Write $j_E^{(2)}$ for the first-stage maps and
$j_E^{(3)}$ for the second-stage maps.  The group $W_3(W_2(U))$
receives the two families
\[
 \alpha_E=j_I^{(3)}\circ j_E^{(2)}\quad(E\in\mathcal P_2),
 \qquad
 \beta_E=j_E^{(3)}\circ j_I^{(2)}\quad(E\in\mathcal P_3).
\]
Postcomposition preserves the first family's identities, and
precomposition preserves the second family's identities.  Their
full-identity maps coincide with the injection
$j_I^{(3)}\circ j_I^{(2)}$.  Notice that the second group algebra is
formed from the abstract group $W_2(U)$; no map between algebras of
different characteristics is involved.

These families satisfy every defining relation of $\widehat W$.
They therefore induce a homomorphism
$\Phi:\widehat W\to W_3(W_2(U))$ for which
\[
 \Phi\circ k_{2,I}=j_I^{(3)}\circ j_I^{(2)}.
\]
The right-hand side is injective, so $k_{2,I}$ is injective as claimed.
\end{proof}

We can now bring the entire finite diagram back into a universal group.

\begin{corollary}\label{prop:self-map}
Let $U$ be a finitely presented group containing an isomorphic copy of
every finitely presented group.  There are an injective endomorphism
$f:U\to U$ and homomorphisms
$h_{\ell,E}:U\to U$ for $\ell\in\{2,3\}$ and
$E\in\mathcal P_\ell$, with $h_{2,I}=h_{3,I}=f$, satisfying all the
orthogonal commutation and sum identities of
Proposition~\ref{prop:finite-envelope}.
\end{corollary}

\begin{proof}
Apply Proposition~\ref{prop:finite-envelope} to $U$.  Since
$\widehat W$ is finitely presented, universality supplies an embedding
$\kappa:\widehat W\to U$.  Set
\[
 h_{\ell,E}=\kappa\circ k_{\ell,E},
 \qquad f=\kappa\circ k_{2,I}=\kappa\circ k_{3,I}.
\]
The map $f$ is injective, and postcomposition by $\kappa$ preserves
every required identity and commutation relation.
\end{proof}

\section{The ascending-torus theorem}
\label{sec:torus}

The finite diagrams constructed above will supply all the higher
finiteness properties at once.  We now state the theorem that makes
this possible and give its proof, including the passage from homology
to a classifying space with finite skeleta.  This ascending-torus
criterion and its homological construction are shared with the
companion~\cite[Theorem~2.2 and Sections~2.4--2.5]{OpenAIOvergroups}.

For a group $U$ and an endomorphism $f:U\to U$, define
\[
 T(U,f)=\langle U,t\mid t^{-1}ut=f(u)\text{ for all }u\in U\rangle.
\]
This presentation includes all the relations of $U$.  When $f$ is
injective, it is the ascending HNN extension associated to $f$.

\Needspace{18\baselineskip}
\begin{theorem}[Ascending-torus theorem]
\label{thm:torus}
Let $U$ be a group and $f:U\to U$ an injective endomorphism.  Suppose
that the following conditions hold.
\begin{enumerate}[label=\textup{(\roman*)}]
\item There are a finitely presented group $P$ and homomorphisms
      $a:U\to P$ and $b:P\to U$ such that $f=ba$.
\item For each $\ell\in\{2,3\}$ there are homomorphisms
      $h_{\ell,E}:U\to U$, indexed by the finite set
      $E\in\mathcal P_\ell$ defined above, with $h_{\ell,I}=f$.
      Whenever $E,E',E+E'\in\mathcal P_\ell$ and
      $EE'=E'E=0$, the subgroups $h_{\ell,E}(U)$ and
      $h_{\ell,E'}(U)$ commute elementwise, and
      \[
       h_{\ell,E+E'}(u)
          =h_{\ell,E}(u)h_{\ell,E'}(u)
          \qquad(u\in U).
      \]
\end{enumerate}
Then $U$ embeds naturally in $T(U,f)$, and $T(U,f)$ is of type
$\Finfty$.
\end{theorem}

The proof has three parts.  First, factoring through a finitely
presented group gives a finite presentation of the ascending torus.
Second, the two idempotent diagrams give uniformly controlled
fillings of ordinary integral cycles after applying an iterate of
$f$.  Finally, these fillings imply vanishing of homology with
products of regular coefficient modules.  A cellular argument then
constructs one classifying space with finitely many cells in every
dimension.  No word-problem hypothesis is used in any of these steps.

\begin{lemma}\label{lem:torus-presentation}
Suppose that $f=ba$ for homomorphisms $a:U\to P$ and $b:P\to U$,
where $P$ is finitely presented.  Then $T(U,f)$ is finitely presented.
This assertion does not require $f$ to be injective.
\end{lemma}

\begin{proof}
Put $g=ab:P\to P$, and write $s$ for the stable letter of $T(P,g)$.
Define maps on the respective generating groups and stable letters by
\[
 \begin{aligned}
 \Phi:T(U,f)&\longrightarrow T(P,g),
 &u&\longmapsto a(u),& t&\longmapsto s,\\
 \Psi:T(P,g)&\longrightarrow T(U,f),
 &p&\longmapsto t b(p)t^{-1},& s&\longmapsto t.
 \end{aligned}
\]
For the first map, $af=ga$ verifies the defining conjugation
relations.  For the second, the image of $s^{-1}ps$ is $b(p)$,
whereas the image of $g(p)$ is
\[
 t b(g(p))t^{-1}=t f(b(p))t^{-1}=b(p).
\]
Thus both are homomorphisms.  Their composites fix the stable letters
and satisfy
\[
 \Psi\Phi(u)=t f(u)t^{-1}=u,
 \qquad
 \Phi\Psi(p)=s g(p)s^{-1}=p.
\]
Consequently $T(U,f)\cong T(P,g)$.  If
$P=\langle y_1,\ldots,y_r\mid R_1,\ldots,R_m\rangle$, choose a word
$w_i$ representing $g(y_i)$ for each $i$.  The finite presentation
\[
 \langle y_1,\ldots,y_r,s\mid
 R_1,\ldots,R_m,\ s^{-1}y_i s=w_i\ (1\le i\le r)\rangle
\]
presents $T(P,g)$.  No injectivity of $g$ is required.
\end{proof}

The remaining task is therefore higher finiteness.  The next section
proves the uniform filling assertion; Section~\ref{sec:products}
uses it to complete the proof of Theorem~\ref{thm:torus}.

\section{Uniform fillings with finitely many simplex types}
\label{sec:fillings}

We now prove the chain-level part of Theorem~\ref{thm:torus}.  Its
purpose is to obtain one iterate of $f$ and one finite set of allowable
simplex types that fill every cycle of a specified dimension and input
bound.  The companion~\cite[Lemma~2.5 and Proposition~2.6]{OpenAIOvergroups}
uses the same homological
device.  It requires neither uniformly short fillings
nor an algorithm for finding them.

\subsection{The homogeneous complex and its bounds}

Let $\mathcal E(U)$ be the simplicial set whose $q$-simplices are all
tuples $(u_0,\ldots,u_q)$ of elements of $U$.  Faces delete entries, and
degeneracies repeat entries.  Write $C_*(U)$ for its normalized integral
chain complex.  Thus its basis in degree $q$ consists of tuples with
no two consecutive entries equal; a tuple with consecutive equal
entries represents zero.  The differential is zero in degree zero,
and, for $q\ge1$, it is
\[
 \partial(u_0,\ldots,u_q)
   =\sum_{i=0}^q(-1)^i(u_0,\ldots,\widehat{u_i},\ldots,u_q).
\]
The augmentation $\epsilon:C_0(U)\to\Z$ sends every vertex to one.
Prepending the identity vertex contracts the augmented complex as a
complex of abelian groups.  Simultaneous left translation gives a free
$U$-action on the normalized basis, so the augmented complex is also
a free resolution of the trivial left $\Z U$-module.  The contraction
need not be $U$-equivariant.

A \emph{bound} is a finite symmetric subset $A\subseteq U$ containing
$1$.  Let $\mathcal E_A(U)$ be the simplicial subset of tuples with
\[
 u_i^{-1}u_j\in A\qquad\text{for all }i,j,
\]
and denote its normalized chains by $C_*^A(U)$.  The condition is
preserved by deleting or repeating vertices, so this is a chain
subcomplex.  A reduced zero-cycle means an element of the kernel of
$\epsilon$; in positive degree, ``reduced cycle'' means an ordinary
cycle.

In each fixed degree, $C_*^A(U)$ has only finitely many basis elements
up to left translation.  Indeed, translating the first vertex to $1$
puts every remaining vertex in $A$.  These orbits are the
\emph{simplex types} of bound $A$.  Notice the distinction: a bound
restricts each simplex's relative vertices, not the locations of its
vertices in $U$, the number of simplices in a chain, or its integer
coefficients.

A homomorphism of groups sends a finite bound to a finite bound.
We say that a chain operator on $C_*^A(U)$ is \emph{controlled} through
specified degrees if there is a finite bound containing every simplex
in all of its output chains in those degrees.  For finitely many
operators and degrees, separate output bounds may be replaced by
their union, with inverses and $1$ adjoined.  For chains in
$\mathcal E(U\times U)$, we will also use separate finite bounds in
the two coordinates.

\subsection{A controlled additivity calculation}

Let $c:C_*(U)\to C_*(U)$ denote the chain map of the constant vertex
map $u\mapsto1$.  Thus $c_0$ sends every vertex to $(1)$, and $c_q=0$
for $q>0$.  The next lemma is an ordinary chain calculation: its
homotopies are not required to commute with the $U$-action.

\Needspace{20\baselineskip}
\begin{lemma}\label{lem:controlled-additivity}
Let $n\ge1$, let $A$ be a bound, and let $f:U\to U$ be an
endomorphism.  Suppose that, for some $k\ge0$, the induced chain map
$F=(f^k)_*:C_*^A(U)\to C_*(U)$ has homomorphisms
\[
 H_q:C_q^A(U)\longrightarrow C_{q+1}(U)
       \qquad(0\le q<n)
\]
whose output chains have a common finite bound and which satisfy
\[
 \partial H_q+H_{q-1}\partial=F_q-c_q
       \qquad(0\le q<n),
\]
where $H_{-1}=0$ and the differential on $C_0$ is zero in this
unaugmented equation.  Let $h_1,h_2:U\to U$ be homomorphisms with
elementwise commuting images, and put $h(u)=h_1(u)h_2(u)$.
Then there is a finite bound $B$ such that every $n$-cycle
$z\in C_n^A(U)$ satisfies
\begin{equation}\label{eq:controlled-additivity}
 [(hf^k)_*z]=[(h_1f^k)_*z]+[(h_2f^k)_*z]
       \quad\text{in }H_n(C_*^B(U)).
\end{equation}
One bound $B$ may be chosen for any fixed finite collection of pairs
$(h_1,h_2)$ satisfying these hypotheses.
\end{lemma}

\begin{proof}
Extend $H$ by zero in all degrees $q\ge n$, and define the chain map
\[
 F'=F-\partial H-H\partial.
\]
Then $F'_0=c_0$ and $F'_q=0$ for $1\le q<n$.  The truncation $H_n=0$
is important: if $z$ is an $n$-cycle, it gives
\begin{equation}\label{eq:truncated-homotopy}
 F'z=Fz.
\end{equation}
Although the two maps need not agree on arbitrary $n$-chains, this
equality is exactly what will be needed.  A common finite output
bound controls $F,F'$ and $H$ in every degree used below; taking
boundaries does not enlarge a bound.

We use the classical Alexander--Whitney and shuffle comparison in the
setting of Eilenberg--Zilber~\cite{EilenbergZilber1953}.  We spell out
the chain maps and their support control: the later product-module
argument requires one finite set of simplex types across all
coordinates.  The Alexander--Whitney diagonal is
\begin{equation}\label{eq:aw-diagonal}
 D(u_0,\ldots,u_q)
   =\sum_{i=0}^q
      (u_0,\ldots,u_i)\otimes(u_i,\ldots,u_q).
\end{equation}
We use the tensor differential
\[
 \partial(a\otimes b)=\partial a\otimes b
                   +(-1)^{\deg a}a\otimes\partial b.
\]
In the boundary of \eqref{eq:aw-diagonal}, terms on the two sides of
an interior cut cancel, and the remaining terms give $D\partial$.
Thus $D$ is a chain map.  It is well defined on normalized chains:
if the input has consecutive equal vertices, every summand has a
degenerate factor.  Both tensor factors retain the input bound.

The shuffle map
\[
 S:C_*(U)\otimes C_*(U)\longrightarrow C_*(U\times U)
\]
sends $(u_0,\ldots,u_p)\otimes(v_0,\ldots,v_q)$ to the signed sum
over monotone lattice paths from $(0,0)$ to $(p,q)$.  A path contributes
the tuple obtained by replacing each of its vertices $(i,j)$ by
$(u_i,v_j)$.  Its sign is $(-1)^N$, where $N$ counts pairs of steps
in which a vertical step precedes a horizontal step.  In its
boundary, deleting a corner pairs two paths that traverse the corner
in opposite orders; these contributions cancel.  The remaining
faces are precisely the shuffle of the tensor differential, so $S$
is a chain map.  If either input is degenerate, every path has a
step with identical endpoint pairs and hence gives a degenerate
output.  The map therefore descends to normalized chains and
preserves the bound in each coordinate.

Write $\Delta$ for the chain map of the diagonal homomorphism
$U\to U\times U$.  There is a chain homotopy $R$ satisfying
\begin{equation}\label{eq:diagonal-homotopy}
 \partial R+R\partial=\Delta-SD
\end{equation}
and preserving the bound in both coordinates.  Here is a direct
construction that also verifies this control.  On vertices take
$R_0=0$, since $\Delta$ and $SD$ agree there.  Suppose $R$ has been
defined in lower degrees, and let $a=(u_0,\ldots,u_q)$ be a
normalized basis simplex of positive degree, with vertex set $V$.
The chain
\[
 \Delta a-SDa-R\partial a
\]
is a cycle: apply the previously established homotopy equation to
$\partial a$.  It lies in the full tuple complex on $V\times V$.
That complex is contracted by prepending any fixed vertex, so
prepending $(u_0,u_0)$ fills the displayed cycle.  Choose this filling
as $R(a)$ and extend linearly on the normalized basis.  Every
coordinate vertex in this construction belongs to $V$, which proves
the claimed bound.  This construction takes place on the normalized
basis; degenerate faces have already been set to zero.

Apply \eqref{eq:diagonal-homotopy} after $F$.  Naturality of $D$
for the vertex map $f^k$ gives a controlled chain homotopy from
$\Delta F$ to $S(F\otimes F)D$.  To replace $F$ by $F'$, use the
degree-one tensor operator $K$ defined on homogeneous tensors by
\begin{equation}\label{eq:tensor-homotopy-signs}
 K(a\otimes b)
     =H(a)\otimes F(b)+(-1)^{\deg a}F'(a)\otimes H(b).
\end{equation}
The tensor differential and the chain-map identities for $F,F'$
give
\begin{align*}
 \partial K+K\partial
   &=(F-F')\otimes F+F'\otimes(F-F')\\
   &=F\otimes F-F'\otimes F'.
\end{align*}
In particular the sign in the second term of
\eqref{eq:tensor-homotopy-signs} cannot be dropped.  The operator
$SKD$ is controlled separately in both coordinates: it uses only
the already controlled maps $F,F',H$ and the bound-preserving maps
$D,S$.  No equivariance of $H$ enters this calculation.
Explicitly, the composite homotopy $L=RF+SKD$ satisfies
\[
 \partial L+L\partial
       =\Delta F-S(F'\otimes F')D.
\]

Let $z$ be an $n$-cycle.  Since $F'$ vanishes in degrees
$1,\ldots,n-1$, only the terms of bidegrees $(0,n)$ and $(n,0)$
survive in $(F'\otimes F')Dz$.  The degree-zero factor always maps
to $(1)$.  Summing the endpoint terms over the whole cycle first, and
then using \eqref{eq:truncated-homotopy}, gives
\[
\begin{aligned}
 (F'\otimes F')Dz
   &=(1)\otimes F'z+F'z\otimes(1)\\
   &=(1)\otimes Fz+Fz\otimes(1).
\end{aligned}
\]
This also covers $n=1$, when there are no intermediate bidegrees.
Shuffling a vertex against a chain simply holds that coordinate
fixed.

Finally, commuting images make
\[
 m:U\times U\longrightarrow U,
       \qquad m(u,v)=h_1(u)h_2(v)
\]
a homomorphism.  If two pairs have coordinate differences in
$B_1$ and $B_2$, their image difference belongs to
$h_1(B_1)h_2(B_2)$.  Indeed, the factors from the two image subgroups
commute, allowing the two differences to be grouped.  This finite
set is symmetric and contains $1$.  Thus applying $m_*$ to all the
preceding homotopies retains finite control and proves
\eqref{eq:controlled-additivity}: the finite chain $m_*Lz$ fills
the difference of its left- and right-hand cycle representatives.
For a finite collection of pairs,
take the union of the finitely many resulting output bounds.
\end{proof}

\subsection{One iterate fills every cycle of a fixed bound}

We now combine the additivity calculation with the two finite
idempotent systems.  The exponent and the output bound in the next
statement are chosen before the input cycle, which is the uniformity
needed when passing later to products of coefficient modules.

\begin{proposition}\label{prop:uniform-fillings}
Assume the hypotheses of Theorem~\ref{thm:torus}.  For every finite
bound $A\subseteq U$ and every integer $n\ge0$, there are an integer
$k\ge0$ and a finite bound $B\subseteq U$ such that, for every
reduced integral $n$-cycle $z\in C_n^A(U)$, there is a finite chain
$y\in C_{n+1}^B(U)$ satisfying
\[
 \partial y=(f^k)_*z.
\]
The chain $y$ may depend on $z$, with no uniform restriction on its
length or coefficients.
\end{proposition}

\begin{proof}
We argue by induction on $n$.  By the finite-presentation
factorization $f=ba$ in Theorem~\ref{thm:torus}, the subgroup $b(P)$
is finitely generated and contains $f(U)$.  Choose a finite symmetric
generating set $S$ for $b(P)$, with $1$ adjoined.  For every $u\in U$,
choose a finite edge path from $1$ to $f(u)$ whose consecutive
differences belong to $S$.  Its chain $H_0(u)$ satisfies
\[
 \partial H_0(u)=(f(u))-(1).
\]
Extend $H_0$ linearly.  For a reduced zero-cycle $z$, the constant
vertex terms cancel, so $\partial H_0(z)=f_*z$.  All these paths use
the same bound $S$, even though their lengths can be arbitrarily
large.  This proves the assertion for $n=0$.

Suppose $n\ge1$ and the assertion holds in all degrees less than
$n$, for every finite input bound.  We first construct the partial
controlled homotopy in Lemma~\ref{lem:controlled-additivity} on
$C_*^A(U)$.  Start with $F=f_*$ and the vertex paths just chosen,
which give its homotopy to the constant map in degree zero.

Suppose, for some $1\le q<n$, that $F=(f^k)_*$ and controlled
operators $H_0,\ldots,H_{q-1}$ have already been constructed with
\[
 \partial H_i+H_{i-1}\partial=F_i-c_i
       \qquad(0\le i<q).
\]
For each normalized basis $q$-simplex $a\in C_q^A(U)$, the next
homotopy equation asks us to fill
\begin{equation}\label{eq:partial-homotopy-obstruction}
 F(a)-H_{q-1}(\partial a).
\end{equation}
This is a $q$-cycle.  Its boundary is zero by the preceding homotopy
equation applied to $\partial a$; the remaining constant-map term
$c_{q-1}(\partial a)$ is zero, also when $q=1$ because the
augmentation of a boundary is zero.  Moreover, all these cycles
have one common finite bound, independent of the basis simplex
$a$: the current $F$ sends $A$ to a finite bound and $H_{q-1}$ is
controlled.

Apply the induction assertion in degree $q$ to this common bound.
It supplies an integer $r\ge0$ and another common finite bound such
that $(f^r)_*$ sends every cycle
\eqref{eq:partial-homotopy-obstruction} to a boundary there.
Postcompose $F$ and the already chosen $H_i$ by $(f^r)_*$.
The existing homotopy equations remain valid, since $f^r(1)=1$.
For each basis simplex choose a finite filling of its newly mapped
cycle and use that filling as $H_q(a)$.  Linear extension defines a
controlled homomorphism $H_q$, and the next homotopy equation now
holds.

There are only finitely many degrees $q<n$, so this process ends
with one power $f^k$ and a controlled homotopy to the constant map
through degrees less than $n$.  For $n=1$ the initial vertex
homotopy already suffices.  None of these choices is required to be
equivariant or computable.  In particular, we do not choose one
filling on each orbit and then translate it.  Uniformity instead
comes from the induction assertion applying at each step to
\emph{every} cycle in a single bounded complex.

Apply Lemma~\ref{lem:controlled-additivity} to every orthogonal
relation in both $\mathcal P_2$ and $\mathcal P_3$.  There are only
finitely many relations, so there is one finite output bound $B$
for all of them; enlarge it if necessary to contain the images of
the input bound under every $h_{\ell,E}f^k$.  For a fixed $n$-cycle $z$,
write
\[
 z_{\ell,E}=[(h_{\ell,E}f^k)_*z]\in H_n(C_*^B(U)).
\]
For each $\ell\in\{2,3\}$, the commuting-image and orthogonal-sum identities,
together with Lemma~\ref{lem:controlled-additivity}, give
\[
 z_{\ell,E+P'}=z_{\ell,E}+z_{\ell,P'}
\]
for every orthogonal triple in that finite system.  Here the three maps
in each relation are evaluated on the common input $f^k(u)$.  The two
full-identity classes agree because both systems have $h_{\ell,I}=f$:
\[
 z_I:=z_{2,I}=z_{3,I}=[(f^{k+1})_*z].
\]
By Lemma~\ref{lem:idempotents}, this single class is annihilated by
both $2$ and $3$.  Hence it is zero, since $z_I=3z_I-2z_I$.
It therefore has a finite filling in $C_*^B(U)$.  The exponent
$k+1$ and the bound $B$ were chosen independently of $z$, proving
the required assertion in degree $n$ and completing the induction.
\end{proof}

Proposition~\ref{prop:uniform-fillings} is a statement about ordinary
integral chains with uniform finite simplex types.  Its homotopies
need not be equivariant, and it asserts no decision procedure for
the word problem of $U$.  The next step uses precisely this form of
control to assemble coordinatewise fillings with product
coefficients.

\section{From controlled fillings to finite skeleta}
\label{sec:products}

Proposition~\ref{prop:uniform-fillings} concerns ordinary integral
chains.  We will apply it coordinate by coordinate to products of
regular modules, keeping the finite set of simplex types common to
all coordinates.  We first explain why the resulting homology
vanishing yields geometric, rather than merely homological,
finiteness.

All coefficient modules in this section are right modules, and all
group homology is integral.  For a group $J$ and a set $\mathcal I$,
write $(\Z J)^{\mathcal I}$ for the direct \emph{product} of copies
of the right regular module, with coordinatewise action.

\subsection{A geometric product-module criterion}

The following argument combines the product-module finiteness method of
Bieri--Eckmann and Brown~\cite{BieriEckmann1974,Brown1975} with the
cellular passage to geometric finiteness.  We give the argument in the
form needed here, including its finite-presentation hypothesis.

\begin{lemma}\label{lem:product-criterion}
Let $J$ be a finitely presented group.  If
\[
 H_n\bigl(J,(\Z J)^{\mathcal I}\bigr)=0
 \qquad\text{for every set }\mathcal I\text{ and every }n\ge2,
\]
then $J$ is of type $\Finfty$.
\end{lemma}

\begin{proof}
We construct nested, $n$-dimensional CW complexes $X_n$, for
$n\ge2$, with fundamental group $J$, finite $n$-skeleton, and
$(n-1)$-connected universal cover.  Each passage from $X_n$ to
$X_{n+1}$ will attach finitely many $(n+1)$-cells.  A finite
presentation complex supplies $X_2$.

Suppose that $X_n$ has been constructed, and put
$Z=\widetilde X_n$.  Thus $Z$ is an $n$-dimensional,
$(n-1)$-connected free $J$-complex.  By successively attaching
$J$-orbits of cells of dimensions at least $n+1$ to kill its
homotopy groups, extend $Z$ to a contractible free $J$-complex $Y$.
There is no finiteness requirement on these auxiliary attachments.
Put $R=\Z J$.  The augmented cellular chain complex $C_*(Y)$ is a
free left $R$-resolution of $\Z$, and $C_q(Y)=C_q(Z)$ is a
finite-rank free $R$-module for $q\le n$.

Let $\mathcal I$ be the set of all cellular $n$-cycles of $Z$, and
let $M=R^{\mathcal I}$.  If $C$ is a finite-rank free left
$R$-module, the coordinate map is an isomorphism
\begin{equation}\label{eq:finite-rank-product-tensor}
 M\otimes_R C\longrightarrow C^{\mathcal I},
 \qquad (m_z)_{z\in\mathcal I}\otimes c
       \longmapsto(m_zc)_{z\in\mathcal I}.
\end{equation}
Using this identification in degrees $n$ and $n-1$, the family
whose $z$-coordinate is the cycle $z$ defines a cycle
$\xi\in M\otimes_R C_n(Y)$.  The hypothesis $H_n(J,M)=0$ gives
\[
 \xi=\partial y
 \qquad\text{for some }y\in M\otimes_R C_{n+1}(Y).
\]
Choose a free $R$-basis of $C_{n+1}(Y)$.  A tensor is a finite sum,
and each chain uses only finitely many basis elements.  Consequently
we can write
\[
 y=\sum_{j=1}^s m_j\otimes c_j
\]
using finitely many basis elements $c_1,\ldots,c_s$, even though
$C_{n+1}(Y)$ need not have finite rank.  Taking the $z$-coordinate
of the boundary equation gives
\[
 z=\sum_{j=1}^s (m_j)_z\,\partial c_j.
\]
Thus the module of cellular $n$-cycles in $Z$ is generated over
$R$ by $\partial c_1,\ldots,\partial c_s$.

Because $Z$ is $(n-1)$-connected and $n\ge2$, the natural Hurewicz map
$\pi_n(Z)\to H_n(Z;\Z)$ is an isomorphism
\cite[Theorem~4.37 and p.~370]{Hatcher2002}.  Naturality with respect to
deck transformations makes it an isomorphism of $R$-modules, using the
canonical change of basepoint in the simply connected space $Z$.
Since $Z$ has no $(n+1)$-cells, its $n$th homology is its module of
cellular $n$-cycles.  Choose maps of $n$-spheres representing the displayed
finite list of generators, and attach their $J$-translates to $Z$.
Downstairs this attaches finitely many $(n+1)$-cells to $X_n$;
upstairs the resulting complex is $n$-connected.  It is the
universal cover of the resulting complex $X_{n+1}$, since cells of
dimension at least three do not change the fundamental group.
This completes the induction.

Let $X=\bigcup_{n\ge2}X_n$.  It has fundamental group $J$ and
finitely many cells in every dimension.  The universal cover of
$X$ is weakly contractible.  It is connected and simply connected, and
for $q\ge2$, cellular approximation \cite[Theorem~4.8]{Hatcher2002}
homotopes any map $S^q\to\widetilde X$ into the $q$-skeleton of
$\widetilde X$, which is already present in $\widetilde X_q$.  That
class is killed in
$\widetilde X_{q+1}$.  A weakly
contractible CW complex is contractible by Whitehead's theorem
\cite[Theorem~4.5]{Hatcher2002}, so $X$ is a $K(J,1)$.
\end{proof}

\subsection{The ascending HNN sequence}

For the rest of the section assume the hypotheses of
Theorem~\ref{thm:torus}, and set $J=T(U,f)$.  We will compute the
homology in Lemma~\ref{lem:product-criterion} using a tree for $J$.
We give the construction explicitly to fix the direction of the
coefficient map.

Form the direct limit
\[
 L=\varinjlim\bigl(U\xrightarrow{f}U\xrightarrow{f}U
          \xrightarrow{f}\cdots\bigr),
\]
and denote its stage maps by $\iota_i:U\to L$, for $i\ge0$.
They satisfy $\iota_i=\iota_{i+1}f$.  Since $f$ is injective,
each $\iota_i$ is injective and
$L=\bigcup_{i\ge0}\iota_i(U)$ is an increasing union.  The map
\[
 \sigma\bigl(\iota_i(u)\bigr)=\iota_{i+1}(u)
\]
is an automorphism: its inverse sends $\iota_i(u)$ to
$\iota_i(f(u))$.  In the semidirect product adjoining a letter
$t$ with $t x t^{-1}=\sigma(x)$ for $x\in L$, we have
\[
 t^{-1}\iota_0(u)t=\iota_0(f(u)).
\]
This gives a homomorphism from $J$ to that semidirect product.
Conversely, the maps $\iota_i(u)\mapsto t^iut^{-i}$ are compatible
with the direct system and with $\sigma$, and give the inverse
homomorphism.  Thus $U$ embeds in $J$.  The same description shows
that every element of $J$ can be written
\begin{equation}\label{eq:ascending-normal-form}
 t^p u t^{-q},\qquad u\in U,\quad p,q\ge0.
\end{equation}
Indeed, an element of $L\rtimes\langle t\rangle$ has the form
$\iota_p(u)t^m$; increasing $p$ if necessary makes $q=p-m$
nonnegative.

Consider the graph whose vertices and oriented edges are both
indexed by the left cosets $jU$, with the edge indexed by $jU$
running from $jU$ to $jtU$.  The terminal vertex is well defined
because
\begin{equation}\label{eq:ascending-commutation}
 ut=tf(u)\qquad(u\in U).
\end{equation}
The homomorphism $J\to\Z$ that kills $U$ and sends $t$ to $1$
defines a height on the vertices.  Each edge increases height by
one, and each vertex has exactly one outgoing edge.
If $j=t^p u t^{-q}$ is as in
\eqref{eq:ascending-normal-form}, the forward path of length $q$
from $jU$ reaches $t^pU$.  Every forward ray therefore meets the
ray $U,tU,t^2U,\ldots$, proving connectedness.  An undirected
cycle would have a vertex of least height with two outgoing edges,
which is impossible.  The graph is a tree.  Left multiplication
gives a $J$-action with one vertex orbit and one edge orbit, and
both the chosen vertex $U$ and the chosen edge from $U$ to $tU$
have stabilizer $U$.

Let $C_*(U)$ be the normalized homogeneous resolution introduced
in the preceding section.  Since $\Z J$ is free as a right
$\Z U$-module, induction preserves its exactness.  We obtain the free left
$\Z J$-resolution
\[
 P_* =\Z J\otimes_{\Z U}C_*(U)
\]
of the permutation module $\Z[J/U]$.  The two endpoint maps on
this resolution are the identity and the map
\[
 E(j\otimes c)=jt\otimes f_*(c).
\]
The second is well defined: for $u\in U$,
\[
 jut\otimes f_*(c)
   =jt f(u)\otimes f_*(c)
   =jt\otimes f_*(uc).
\]
It is a left $\Z J$-linear chain map, lifting the terminal-vertex
map $jU\mapsto jtU$.

The mapping cone of $1-E:P_*\to P_*$, with its natural
augmentation, is a free left $\Z J$-resolution of $\Z$.
To see exactness, first take the homology of $P_*$: it is
concentrated in degree zero and equal to $\Z[J/U]$.  The
remaining two-term complex is the augmented cellular chain
complex of the tree, which is exact.  Equivalently, the long exact
sequence of the cone identifies its homology with that of this
two-term complex.

For an arbitrary right $\Z J$-module $M$, tensoring this free
resolution with $M$ yields the long exact sequence
\begin{equation}\label{eq:ascending-hnn-homology}
 \begin{split}
 H_n(U,M)&\xrightarrow{\,1-F\,}H_n(U,M)
      \longrightarrow H_n(J,M)\\
      &\longrightarrow H_{n-1}(U,M)
          \xrightarrow{\,1-F\,}H_{n-1}(U,M),
 \end{split}
\end{equation}
where $M$ is restricted to $U$ in the outer terms.  On the chain
complex $M\otimes_{\Z U}C_*(U)$, the map $F$ is
\begin{equation}\label{eq:coefficient-shift}
 F(m\otimes c)=mt\otimes f_*(c).
\end{equation}
In particular the coefficient multiplier is $t$, not $t^{-1}$.

\subsection{Products of regular modules}

Fix any set $\mathcal I$ and take $M=(\Z J)^{\mathcal I}$.
We claim that $F$ induces a locally nilpotent endomorphism of
$H_i(U,M)$ for every $i>0$: each individual homology class is
killed by some power of $F$.

Recall that a bound is a finite symmetric subset $A\subseteq U$
containing $1$, and $C_q^A(U)$ is generated by the normalized
tuples $(u_0,\ldots,u_q)$ with $u_i^{-1}u_j\in A$ for all $i,j$.
Let $\mathcal S_q(A)$ be their finite set of types up to left
translation.  The first vertex can be translated to $1$, after
which all other vertices lie in $A$, proving finiteness.  As
$C_q^A(U)$ is free over $\Z U$ on these types, there are natural
identifications
\begin{equation}\label{eq:bounded-product-chains}
 \begin{split}
 M\otimes_{\Z U}C_q^A(U)
 &\cong\bigoplus_{s\in\mathcal S_q(A)}M\\
 &\cong\prod_{\alpha\in\mathcal I}
      \left(\bigoplus_{s\in\mathcal S_q(A)}\Z J\right).
 \end{split}
\end{equation}
The second equality interchanges a product with a \emph{finite}
direct sum; no analogous interchange with an infinite direct sum
is being used.

Here is a concrete interpretation of
\eqref{eq:bounded-product-chains}.  The free abelian basis of
$\Z J\otimes_{\Z U}C_q^A(U)$ consists of normalized tuples with
$u_0,\ldots,u_q\in U$ of the form
\[
 (ju_0,\ldots,ju_q),\qquad j\in J,
 \quad u_a^{-1}u_b\in A\text{ for all }a,b.
\]
Each tuple lies in one left coset of $U$.  Expanding the group-ring
coefficients, a tensor chain therefore gives, in each coordinate
$\alpha$, a finite ordinary integral chain in such coset tuples.
There is one bound $A$ for every coordinate.  Conversely, any
family of coordinatewise finite chains with a common bound
defines a tensor chain by \eqref{eq:bounded-product-chains}.
Their lengths, coefficients, and group-ring supports need not
admit any bound independent of $\alpha$.

Every chain in $M\otimes_{\Z U}C_q(U)$ uses finitely many
simplex types and consequently belongs to a bounded subcomplex
of this form.  Indeed, take the union of the finitely many
relative vertex differences in its tensor factors, then add
their inverses and $1$.  The inclusion of this bounded complex
is injective degree by degree, since its types are a subset of
the free basis of $C_q(U)$.

Now let $z\in M\otimes_{\Z U}C_i(U)$ be a cycle, where $i>0$,
and choose one input bound $A$ for it.  In each coordinate
$\alpha$, split its ordinary chain into its finitely many
left-coset components.  Each component is a cycle, since the
boundary preserves the coset and the coset chain groups form a
direct sum.  Choose an origin $j$ in each nonzero component's
coset.  Left translation by $j^{-1}$ identifies the component
with an ordinary cycle in $C_i^A(U)$.

Proposition~\ref{prop:uniform-fillings} supplies an exponent
$k\ge0$ and a finite output bound $B$ that work for
\emph{every} such cycle, in every coordinate and every coset.
By \eqref{eq:coefficient-shift}, the map $F^k$ changes the
origin $j$ to $jt^k$ and applies $f^k$ to the relative vertices.
Indeed, on tuples the single-step formula is simply
\[
 (ju_0,\ldots,ju_i)
    \longmapsto(jt f(u_0),\ldots,jt f(u_i))
    =(ju_0t,\ldots,ju_it).
\]
This also verifies that the formula is independent of the
chosen origin.

For each component, choose an ordinary finite filling of its
relative image under $f^k$ in $C_{i+1}^B(U)$, and translate
this filling by $jt^k$.  Sum the fillings within each coordinate.
The result is still finite in that coordinate because there were
only finitely many nonzero coset components.  Different input
cosets may have the same output coset; in that case their
fillings are simply added.  Across all coordinates the bound
is still the single finite set $B$.
Equation~\eqref{eq:bounded-product-chains} therefore assembles
these fillings into a chain
\[
 y\in M\otimes_{\Z U}C_{i+1}^B(U).
\]
The boundary and $F^kz$ lie in the bounded complex for
$D=B\cup f^k(A)$, a finite symmetric set containing $1$.  They agree
in every coordinate by construction.  The coordinate identification
\eqref{eq:bounded-product-chains} for $D$ is injective, so
$\partial y=F^kz$ in the tensor complex.
Thus $F^k[z]=0$, proving the claimed local nilpotence.  This
argument uses ordinary fillings chosen separately in each coset;
it requires no equivariance of the homotopies or fillings in
Proposition~\ref{prop:uniform-fillings}.

\subsection{Completion of the torus theorem}

If an endomorphism $F$ of an abelian group is locally nilpotent,
then $1-F$ is invertible.  On an element $x$, its inverse is
\[
 x+Fx+F^2x+\cdots,
\]
where only finitely many terms are nonzero.  The sum defines a
homomorphism, and telescoping verifies both inverse identities.
It follows that the two maps $1-F$ adjacent to $H_n(J,M)$ in
\eqref{eq:ascending-hnn-homology} are isomorphisms for $n\ge2$.
Exactness gives
\[
 H_n\bigl(J,(\Z J)^{\mathcal I}\bigr)=0
 \qquad(n\ge2).
\]
Only positive-degree local nilpotence was needed: no assertion
about $H_0(U,M)$ enters this conclusion.

By Lemma~\ref{lem:torus-presentation}, $J$ is finitely presented.
The set $\mathcal I$ was arbitrary, so
Lemma~\ref{lem:product-criterion} proves that $J$ is of type
$\Finfty$.  The direct-limit construction already proved that
$U$ embeds in $J$.  This completes the proof of
Theorem~\ref{thm:torus}.

\section{Proof of the embedding theorem}\label{sec:assembly}

We now combine the finite-presentation construction with the geometric
criterion.

\begin{proof}[Proof of Theorem~\ref{thm:main}]
Let $U$ be the universal finitely presented group supplied by
Lemma~\ref{lem:universal}. Corollary~\ref{prop:self-map} gives an
injective endomorphism $f:U\to U$ and, for each $\ell\in\{2,3\}$,
homomorphisms $h_{\ell,E}:U\to U$ indexed by $E\in\mathcal P_\ell$.
Their full-identity maps are both $f$, and they satisfy all the
elementwise commutation and orthogonal-sum identities in
Theorem~\ref{thm:torus}.

The factorization hypothesis of that theorem is automatic here: take
$P=U$, $a=\operatorname{id}_U$, and $b=f$. Hence
$H=T(U,f)$ has type $\Finfty$. Since $f$ is injective, the natural map
$U\to H$ is injective, as established in Section~\ref{sec:products}.
For every finitely presented group $G$, choose an embedding $G\to U$.
Its composite with $U\to H$ is the required embedding. The construction
of $H$ is independent of this last choice and of $G$.
\end{proof}

The construction isolates a useful distinction: finite presentation is
needed for the universal group and its diagram envelope, while the
faithful models in tensor product algebras are used only to prove injectivity.
The higher-dimensional finiteness is obtained afterwards, from the
two finite systems and one ascending HNN extension.

\begin{thebibliography}{9}

\bibitem{BieriEckmann1974}
Robert Bieri and Beno Eckmann,
\emph{Finiteness properties of duality groups},
Commentarii Mathematici Helvetici \textbf{49} (1974), 74--83.
\href{https://doi.org/10.1007/BF02566720}{doi:10.1007/BF02566720}.

\bibitem{Brown1975}
Kenneth S.\ Brown,
\emph{Homological criteria for finiteness},
Commentarii Mathematici Helvetici \textbf{50} (1975), 129--135.
\href{https://doi.org/10.1007/BF02565740}{doi:10.1007/BF02565740}.

\bibitem{EilenbergZilber1953}
Samuel Eilenberg and J.\ A.\ Zilber,
\emph{On products of complexes},
American Journal of Mathematics \textbf{75} (1953), no.~1, 200--204.
\href{https://doi.org/10.2307/2372629}{doi:10.2307/2372629}.

\bibitem{FFZ2026}
Francesco Fournier-Facio and Matthew C.\ B.\ Zaremsky,
\emph{Finiteness properties and Higman's rope trick},
preprint, 2026,
\href{https://arxiv.org/abs/2607.21727v1}{arXiv:2607.21727v1}.

\bibitem{Hatcher2002}
Allen Hatcher,
\emph{Algebraic Topology},
Cambridge University Press, 2002.
\url{https://pi.math.cornell.edu/~hatcher/AT/ATpage.html}.

\bibitem{Higman1961}
Graham Higman,
\emph{Subgroups of finitely presented groups},
Proc. Roy. Soc. London Ser. A \textbf{262} (1961), 455--475.
\href{https://doi.org/10.1098/rspa.1961.0132}{doi:10.1098/rspa.1961.0132}.

\bibitem{HNN1949}
Graham Higman, B.\ H.\ Neumann and Hanna Neumann,
\emph{Embedding theorems for groups},
J. London Math. Soc. \textbf{24} (1949), no.~4, 247--254.
\href{https://doi.org/10.1112/jlms/s1-24.4.247}{doi:10.1112/jlms/s1-24.4.247}.

\bibitem{Leary2018}
Ian J.\ Leary,
\emph{Subgroups of almost finitely presented groups},
Math. Ann. \textbf{372} (2018), 1383--1391.
\href{https://doi.org/10.1007/s00208-018-1689-5}{doi:10.1007/s00208-018-1689-5}.

\bibitem{OpenAIOvergroups}
OpenAI,
\emph{Simple $F_\infty$ overgroups of groups with decidable word problem},
OpenAI Math Release preprint
\href{https://github.com/openai/math/blob/main/preprints/Simple-F-infinity-overgroups-of-groups-with-decidable-word-problem-September-23-2026/paper.pdf}{OAI:Simple-F-infinity-overgroups-of-groups-with-decidable-word-problem-September-23-2026},
2026.

\end{thebibliography}

\end{document}
